\section*{Erd\H{o}s Problem 705: large girth and unit distances}

\paragraph{Statement and status of this reconstruction.}
The question asks whether sufficiently large girth forces a
finite planar unit-distance graph to be \(3\)-colourable,
with an edge required \emph{if and only if} the distance is one.
The problem page records O'Donnell's negative resolution:
\url{https://www.erdosproblems.com/705}.
We reconstruct the graph-theoretic part of his construction
in detail. We use his published embedding theorem for its
stated, weaker edge-preserving conclusion, and explicitly
leave the faithful-embedding step unresolved in this attempt.
This distinction is about the scope verified here, not a
claim that the recorded problem is still open.

\paragraph{Sources and the geometric qualification.}
The primary source inspected is P. O'Donnell,
\emph{High girth unit-distance graphs}, Rutgers Ph.D. thesis
(1999), Sections 1.2 and 1.7, Theorems 25--28:
\url{https://jayyhk.github.io/papers/odonnell1999.pdf}.
Its definition of a proper unit-distance embedding requires
distinct points and unit length for every edge.
That definition alone does not require nonedges to have
length different from one. The discussion at
\url{https://www.erdosproblems.com/forum/thread/705}
accepts the negative resolution and also mentions this
faithfulness issue.
We do not infer the missing converse just from the word
proper or from the statement of Theorem 28.

\paragraph{An explicit graph construction from a hypergraph.}
Choose an odd integer \(k\geq5\) and an integer \(g>k/3\),
with \(g\geq3\).
Use the established Erd\H{o}s--Hajnal theorem giving a finite
\(k\)-uniform hypergraph \(\mathcal H\) of Berge girth at
least \(g\) which is not weakly \(3\)-colourable.
Weak colouring means that no hyperedge is monochromatic.
This existence theorem is an external combinatorial input,
stated as Theorem 25 in the cited thesis.

Start with one independent foundation vertex for each vertex
of \(\mathcal H\). For every hyperedge
\(E=\{v_1,\ldots,v_k\}\), introduce a fresh cycle
\(u_{E,1},\ldots,u_{E,k}\) of length \(k\), and add the matching
edges \(v_i u_{E,i}\).
Cycles belonging to different hyperedges have disjoint new
vertices; foundation vertices may belong to several hyperedges.
Call the resulting graph \(G\).

\paragraph{Its chromatic number is exactly four.}
In any putative \(3\)-colouring of \(G\), the foundation
vertices give a \(3\)-colouring of \(\mathcal H\). Some
hyperedge \(E\) is monochromatic, say in colour one.
Every vertex of its attached cycle is adjacent to a
foundation vertex of colour one, so that odd cycle would
have to use just colours two and three, a contradiction.
Hence \(\chi(G)\geq4\).
Conversely, give every foundation vertex one colour and
properly colour every attached odd cycle with three
other colours. This proves \(\chi(G)=4\).

\paragraph{Its girth is exactly \(k\).}
An attached cycle shows that the girth is at most \(k\).
Consider any other simple cycle \(C\).
If it has no foundation vertex, it lies wholly in one
attached cycle and hence has length \(k\).
Otherwise, between consecutive foundation vertices on
\(C\) there is a path inside a single attached cycle,
including the two matching edges. This path has at least
three edges.

Project these successive paths to the incidence graph of
\(\mathcal H\), replacing each by
foundation vertex--hyperedge--foundation vertex.
Consecutive hyperedges in this projected closed walk are
different: at a fixed foundation vertex there is only one
matching edge into each particular attached cycle, so
using the same hyperedge immediately twice would repeat
that matching edge in the simple cycle \(C\).
The foundation vertices on \(C\) are also distinct.
Thus the projected closed walk has no immediate reversal.
It contains a cycle in the incidence graph, whose length
is at least \(2g\) by the Berge-girth assumption.
If \(C\) has \(t\) foundation vertices, its projection
has length \(2t\), so \(t\geq g\) and
\(|C|\geq3t\geq3g>k\).
This proves the girth assertion without assuming that
all projected hyperedges are distinct.

\paragraph{What the external embedding result supplies.}
Choose \(\mathcal H\) vertex-minimal among hypergraphs with
these properties. Deleting any one foundation vertex makes
its remaining hypergraph weakly \(3\)-colourable.
O'Donnell's embedding construction uses this colouring and
four small regions of the plane, then attaches the odd
cycles and removes coincidences. The geometric lemmas and
Theorem 28 give distinct planar points realizing every
specified edge at distance one. This geometric existence
is an external input, not established by the colouring
and girth arguments above.

\paragraph{Remaining gap for the displayed stronger convention.}
To turn this into a complete proof of the stated
if-and-only-if version, one must verify a realization in
which every unspecified vertex pair has distance different
from one, or otherwise control all additional unit edges.
Simply adding every missing unit edge can introduce short
cycles and destroys the proved girth bound.
The source's removal of coincident vertices alone does not
address that issue. Accordingly the faithful geometric
step remains unresolved in this reconstruction, even
though the abstract construction and its one-way
unit-distance realization have the properties proved above.

\paragraph{Completion estimate.}
\textbf{COMPLETION: 65\%}
